% Сборка: xelatex lab01.tex
\documentclass[a4paper,17pt]{extarticle}

\usepackage{fontspec}
\usepackage{polyglossia}
\setmainlanguage{russian}
\setotherlanguage{english}

% Рубленые шрифты лучше читаются при слабом зрении.
\setmainfont{Liberation Sans}
\setsansfont{Liberation Sans}
\setmonofont{Liberation Mono}[Scale=MatchLowercase]

\usepackage{amsmath}
\usepackage{geometry}
\geometry{a4paper, margin=2.2cm}
\usepackage{setspace}
\setstretch{1.35}
\usepackage{enumitem}
\usepackage{listings}
\usepackage{xcolor}

% Строго чёрный текст на белом фоне, без цветной подсветки.
\pagecolor{white}
\color{black}

\setlist[enumerate]{leftmargin=*, itemsep=1.1em, topsep=0.6em}
\raggedright
\hyphenpenalty=10000
\exhyphenpenalty=10000
\setlength{\parindent}{0pt}
\setlength{\parskip}{0.6em}

% Чёрно-белое оформление кода: ключевые слова жирным, комментарии курсивом,
% без цвета и без серой заливки, толстая чёрная рамка.
\lstdefinestyle{bw}{
  language=Python,
  basicstyle=\ttfamily\normalsize\color{black},
  literate={->}{{-{>}}}2,
  keywordstyle=\bfseries\color{black},
  commentstyle=\itshape\color{black},
  stringstyle=\color{black},
  identifierstyle=\color{black},
  numbers=none,
  frame=single,
  rulecolor=\color{black},
  framerule=1.2pt,
  backgroundcolor=\color{white},
  xleftmargin=0.6em,
  xrightmargin=0.6em,
  framexleftmargin=0.6em,
  framexrightmargin=0.6em,
  framextopmargin=0.4em,
  framexbottommargin=0.4em,
  aboveskip=0.8em,
  belowskip=0.8em,
  columns=fullflexible,
  keepspaces=true,
  showstringspaces=false,
  breaklines=true,
  upquote=true,
}
\lstset{style=bw}

\newcommand{\pts}[1]{\textbf{(#1)}}

\begin{document}

\begin{center}
  {\Large\bfseries Лабораторная работа 1\par}
  {\large Точечные оценки параметров распределения\par}
\end{center}

\textbf{Как сдавать работу.}
\begin{enumerate}[itemsep=0.3em]
  \item Сохраните копию ноутбука на свой Диск («Файл $\to$ Сохранить копию на Диске»).
  \item Решите задания в своей копии.
  \item Откройте доступ: «Поделиться $\to$ Все, у кого есть ссылка $\to$ Читатель».
  \item Пришлите преподавателю ссылку на свою копию.
\end{enumerate}

Всего за работу можно получить 10 баллов.

\section*{Теория}

\begin{enumerate}
  \item \pts{1 балл} Дайте определение точечной оценки параметров распределения по случайной выборке.

  \item \pts{1 балл} Что такое несмещенная оценка? Что такое состоятельная оценка?

  \item \pts{1 балл} Напишите формулу плотности одномерного нормального (гауссовского) распределения. Для его параметров $\mu$ и $\sigma^2$ (оба неизвестны) напишите формулы несмещенных оценок по выборке $X=\left\{x_n\right\}_{n=1}^N$.

  \item \pts{2 балла} Докажите, что приведенные в пункте 3 оценки $\hat{\mu}_N$ и $\hat{\sigma}_N^2$ являются несмещенными и состоятельными.
\end{enumerate}

\section*{Практика}

\begin{enumerate}
  \item \pts{1 балл} Выберите какие-нибудь значения математического ожидания $\mu$ и дисперсии $\sigma^2$ нормального распределения и с помощью генератора \lstinline|np.random.default_rng(seed)| и его метода \lstinline|rng.normal| сгенерируйте выборку объема $N=1000$. Обратите внимание, что параметр \lstinline|scale| в \lstinline|rng.normal|~--- это стандартное отклонение $\sigma$, а не дисперсия $\sigma^2$. Постройте на одном графике плотность нормального распределения и гистограмму выборки (с нормировкой \lstinline|density=True|) и убедитесь, что гистограмма приближает плотность.

  \item \pts{1 балл} Реализуйте функцию, вычисляющую несмещенную оценку математического ожидания $\hat{\mu}_N$ по случайной выборке $X=\left\{x_n\right\}_{n=1}^N$ объема $N$ нормального распределения:
\begin{lstlisting}
def hat_mu(X: np.ndarray) -> float:
    # ваш код
    return mu
\end{lstlisting}

  \item \pts{1 балл} Реализуйте функцию, вычисляющую несмещенную оценку дисперсии $\hat{\sigma}_N^2$ по случайной выборке $X=\left\{x_n\right\}_{n=1}^N$ объема $N$ нормального распределения (математическое ожидание считайте неизвестным):
\begin{lstlisting}
def hat_var(X: np.ndarray) -> float:
    # ваш код
    return var
\end{lstlisting}
  Проверьте свои функции, сравнив результаты с \lstinline|np.mean(X)| и \lstinline|np.var(X, ddof=1)|.

  \item \pts{2 балла} Для выбранных в задаче 1 значений математического ожидания $\mu$ и дисперсии $\sigma^2$ постройте графики зависимостей оценок $\hat{\mu}_N$ и $\hat{\sigma}_N^2$ от объема выборки $N$ (например, по первым $N$ элементам одной большой выборки) и отметьте на них истинные значения параметров. Убедитесь, что при увеличении $N$ оценки приближаются к истинным значениям параметров.
\end{enumerate}

\end{document}
